\documentclass[11pt,letterpaper]{article}

\usepackage[utf8]{inputenc}
\usepackage[spanish]{babel}
\usepackage[margin=2.5cm]{geometry}
\usepackage{amsmath,amssymb}
\usepackage{enumitem}
\usepackage{titlesec}
\usepackage{array}
\usepackage{hyperref}
\usepackage{xcolor}

\definecolor{unicaucaverde}{RGB}{0,90,60}
\hypersetup{colorlinks=true,linkcolor=unicaucaverde,urlcolor=unicaucaverde}

\titleformat{\section}{\normalfont\Large\bfseries\color{unicaucaverde}}{\thesection}{1em}{}
\titleformat{\subsection}{\normalfont\large\bfseries}{\thesubsection}{1em}{}

\newcommand{\vv}[1]{\mathbf{#1}}

\setlength{\parindent}{0pt}
\setlength{\parskip}{6pt}

\title{\textcolor{unicaucaverde}{\Large Guía de Aprendizaje — Semana 4}\\[4pt]
  \large Matrices, matriz inversa y transpuesta}
\author{Álgebra Lineal para Ingeniería --- Universidad del Cauca}
\date{2026}

\begin{document}
\maketitle

\begin{center}
\begin{tabular}{|>{\bfseries}p{4cm}|p{10cm}|}
\hline
Semana & 4 de 16 \\ \hline
Unidad & 2 --- Vectores y matrices (Grossman, cap. 2) \\ \hline
Subtemas & 2.3 (matrices y sistemas), 2.4 (inversa), 2.5 (transpuesta) \\
\hline
Horas de clase & 4 \\ \hline
Resultado de aprendizaje & RA2 (parte 2 de 3; ver \texttt{guia1/guia.tex}) \\ \hline
\end{tabular}
\end{center}

\section*{Relevancia para ingeniería}

En ingeniería, la matriz de coeficientes $A$ representa la estructura de un sistema dinámico, red
o circuito, mientras que el vector de entrada o excitación $\vv{b}$ produce la respuesta $\vv{x}$.
Resolver el sistema mediante la matriz inversa $\vv{x}=A^{-1}\vv{b}$ permite calcular de forma directa
la respuesta frente a múltiples estados de entrada sin necesidad de rehacer la eliminación gaussiana
cada vez, lo cual es vital en control automático, análisis estructural y procesamiento de señales.
Por su parte, la operación transpuesta $A^T$ aparece de forma natural al analizar la energía de un
sistema, en mínimos cuadrados, en la transformación de marcos de referencia (gráfica por computador y
robótica) y en el estudio de matrices simétricas (que caracterizan tensores de esfuerzo y matrices de
covarianza).

\section{Objetivo de la semana}

Expresar sistemas de ecuaciones lineales en forma matricial $A\vv{x}=\vv{b}$, calcular la inversa de
matrices cuadradas mediante el método de Gauss-Jordan y fórmulas explícitas, resolver sistemas usando
la matriz inversa, y aplicar las propiedades algebraicas de la transpuesta y las matrices simétricas.

\section{Resultados de aprendizaje de la semana}

\begin{itemize}[leftmargin=1.2cm]
  \item[\textbf{RA2.3.}] Representar un sistema de $m$ ecuaciones con $n$ incógnitas en forma
        matricial $A\vv{x}=\vv{b}$ e interpretar el papel del rango y la consistencia en términos
        de operaciones matriciales.
  \item[\textbf{RA2.4.}] Determinar si una matriz cuadrada es invertible (no singular), calcular su
        inversa $A^{-1}$ mediante Gauss-Jordan $[A|I]\to [I|A^{-1}]$ o la fórmula $2\times 2$, y
        resolver sistemas de ecuaciones mediante $\vv{x}=A^{-1}\vv{b}$.
  \item[\textbf{RA2.5.}] Calcular la transpuesta $A^T$ de cualquier matriz, identificar matrices
        simétricas ($A^T=A$) y antisimétricas ($A^T=-A$), y aplicar las propiedades algebraicas
        de la transpuesta, incluyendo $(AB)^T = B^T A^T$ y $(A^T)^{-1} = (A^{-1})^T$.
\end{itemize}

\section{Contenidos}

\begin{itemize}[leftmargin=1.2cm]
  \item \textbf{2.3.} Matrices y sistemas de ecuaciones lineales: notación matricial $A\vv{x}=\vv{b}$,
        matriz de coeficientes y matriz aumentada.
  \item \textbf{2.4.} Inversa de una matriz cuadrada: definición, condición de invertibilidad,
        algoritmo de Gauss-Jordan $[A|I]\to [I|A^{-1}]$, inversa de matriz $2\times 2$, propiedades
        y solución de sistemas $A\vv{x}=\vv{b}$.
  \item \textbf{2.5.} Transpuesta de una matriz: definición, matrices simétricas y antisimétricas,
        propiedades algebraicas $(AB)^T=B^T A^T$ y relación entre transposición e inversión.
\end{itemize}

\section{Plan de la sesión (4 horas)}

\begin{enumerate}[leftmargin=1.2cm]
  \item[\textbf{Hora 1.}] Matrices y sistemas de ecuaciones lineales: forma matricial $A\vv{x}=\vv{b}$,
        interpretación de ecuaciones como producto de vectores fila por vector incógnita.
  \item[\textbf{Hora 2.}] Concepto de matriz inversa $A^{-1}$; matrices singulares y no singulares;
        fórmula para inversa de matriz $2\times 2$; algoritmo $[A|I]\to [I|A^{-1}]$ mediante
        operaciones elementales de fila.
  \item[\textbf{Hora 3.}] Propiedades de la matriz inversa: $(AB)^{-1}=B^{-1}A^{-1}$, $(c A)^{-1}=c^{-1}A^{-1}$;
        solución directa de sistemas de ecuaciones $\vv{x}=A^{-1}\vv{b}$ cuando $A$ es invertible.
  \item[\textbf{Hora 4.}] Transpuesta de una matriz $A^T$; propiedades del producto $(AB)^T=B^TA^T$;
        matrices simétricas y antisimétricas; taller práctico con los 10 ejercicios resueltos de \texttt{notas.tex}.
\end{enumerate}

\section{Actividades}

\begin{itemize}[leftmargin=1.2cm]
  \item \textbf{En clase}: taller en grupo durante la Hora 4 resolviendo ejercicios de inversión y
        transposición de matrices.
  \item \textbf{Trabajo autónomo}: estudio del material de \texttt{notas.tex} y resolución de la lista
        de ejercicios previa a la Semana 5.
\end{itemize}

\section{Material de apoyo}

\begin{itemize}[leftmargin=1.2cm]
  \item \texttt{notas.tex} --- teoría detallada de 2.3, 2.4 y 2.5, con demostraciones, ejemplos y
        10 ejercicios resueltos paso a paso.
  \item \texttt{diapositivas.tex} --- presentación sintética para la clase magistral.
  \item \texttt{material-web/} --- plataforma interactiva con calculadoras web de matriz inversa y
        transpuesta.
\end{itemize}

\section{Evaluación de la semana}

Control formativo corto al inicio de la Semana 5 (10 minutos): dada una matriz $2\times 2$ o $3\times 3$,
calcular su inversa (si existe), resolver un sistema $A\vv{x}=\vv{b}$ usando $A^{-1}$, y hallar $A^T$.

\section{Bibliografía de la semana}

\begin{enumerate}[leftmargin=1.2cm]
  \item Grossman, S. I., Flores Godoy, J. J. \emph{Álgebra lineal}. 7a/8a ed., McGraw-Hill,
        cap. 2, §2.3--2.5.
  \item Lay, D. C., Lay, S. R., McDonald, J. J. \emph{Álgebra lineal y sus aplicaciones}. Pearson, cap. 2.
\end{enumerate}

\end{document}
